% Options for packages loaded elsewhere
\PassOptionsToPackage{unicode}{hyperref}
\PassOptionsToPackage{hyphens}{url}
\PassOptionsToPackage{dvipsnames,svgnames,x11names}{xcolor}
%
\documentclass[
  11pt,
]{article}
\usepackage{amsmath,amssymb}
\usepackage{iftex}
\ifPDFTeX
  \usepackage[T1]{fontenc}
  \usepackage[utf8]{inputenc}
  \usepackage{textcomp} % provide euro and other symbols
\else % if luatex or xetex
  \usepackage{unicode-math} % this also loads fontspec
  \defaultfontfeatures{Scale=MatchLowercase}
  \defaultfontfeatures[\rmfamily]{Ligatures=TeX,Scale=1}
\fi
\usepackage{lmodern}
\ifPDFTeX\else
  % xetex/luatex font selection
\fi
% Use upquote if available, for straight quotes in verbatim environments
\IfFileExists{upquote.sty}{\usepackage{upquote}}{}
\IfFileExists{microtype.sty}{% use microtype if available
  \usepackage[]{microtype}
  \UseMicrotypeSet[protrusion]{basicmath} % disable protrusion for tt fonts
}{}
\makeatletter
\@ifundefined{KOMAClassName}{% if non-KOMA class
  \IfFileExists{parskip.sty}{%
    \usepackage{parskip}
  }{% else
    \setlength{\parindent}{0pt}
    \setlength{\parskip}{6pt plus 2pt minus 1pt}}
}{% if KOMA class
  \KOMAoptions{parskip=half}}
\makeatother
\usepackage{xcolor}
\usepackage[margin=1in]{geometry}
\usepackage{longtable,booktabs,array}
\usepackage{calc} % for calculating minipage widths
% Correct order of tables after \paragraph or \subparagraph
\usepackage{etoolbox}
\makeatletter
\patchcmd\longtable{\par}{\if@noskipsec\mbox{}\fi\par}{}{}
\makeatother
% Allow footnotes in longtable head/foot
\IfFileExists{footnotehyper.sty}{\usepackage{footnotehyper}}{\usepackage{footnote}}
\makesavenoteenv{longtable}
\setlength{\emergencystretch}{3em} % prevent overfull lines
\providecommand{\tightlist}{%
  \setlength{\itemsep}{0pt}\setlength{\parskip}{0pt}}
\setcounter{secnumdepth}{-\maxdimen} % remove section numbering
\ifLuaTeX
  \usepackage{selnolig}  % disable illegal ligatures
\fi
\usepackage{bookmark}
\IfFileExists{xurl.sty}{\usepackage{xurl}}{} % add URL line breaks if available
\urlstyle{same}
\hypersetup{
  pdftitle={Three sign changes cannot improve the nine-weight Conway--Guy witness},
  pdfauthor={Aafi Khan Malek (Independent Researcher, Ahmedabad, India)},
  colorlinks=true,
  linkcolor={Maroon},
  filecolor={Maroon},
  citecolor={Blue},
  urlcolor={Blue},
  pdfcreator={LaTeX via pandoc}}

\title{Three sign changes cannot improve the nine-weight Conway--Guy
witness}
\usepackage{etoolbox}
\makeatletter
\providecommand{\subtitle}[1]{% add subtitle to \maketitle
  \apptocmd{\@title}{\par {\large #1 \par}}{}{}
}
\makeatother
\subtitle{An exact local exclusion by positive dual certificates}
\author{Aafi Khan Malek (Independent Researcher, Ahmedabad, India)}
\date{November 2026}

\begin{document}
\maketitle

\section{The nine-weight Conway--Guy sign-radius-three
obstruction}\label{the-nine-weight-conwayguy-sign-radius-three-obstruction}

\textbf{Theorem.} Let \(x_1<\cdots<x_9\) be positive real weights with
pairwise unit-separated subset sums. If their signed-relation chamber
differs in at most three signs from the Conway--Guy witness
\(a=(77,117,137,148,154,157,159,160,161)\), then \(x_9\ge161\). A sign
here is one ternary word up to total reversal, not one literal subset
comparison or a Euclidean distance. No assertion is made about distance
four or more.

The exact radius-two package is included unmodified under
prior-radius2/. There, the 9,841 oriented words, all 12 tight relations,
rational single-wall crossings and objective certificates are
independently checked. The present proof starts with
\href{https://arxiv.org/abs/2611.17360}{Notik's identity}
\(x_9=\sum_{i=1}^9y_iR_i\), with \[
(y_1,\ldots,y_9)=(1,1,2,4,8,15,25,48,57).
\] Each \(R_i\) is a ternary word, CG-positive with value 1 at \(a\);
the explicit table is in prior-radius2/README.md. A positive dual is an
identity expressing \(x_9\) as a nonnegative combination of the signed
rows enforced in the chamber. Each such row is at least 1, so its
coefficient sum is a valid lower bound. We give a complete support
reduction and check every resulting nontrivial combination by exact
integer arithmetic. No optimizer is needed for the proof.

\subsection{Notation and substitution
rules}\label{notation-and-substitution-rules}

\(R_8=x_8-x_7\) and \(R_9=x_9-x_8\) are always positive in a sorted
chamber. Thus only \(R_1,\ldots,R_7\) can be among the changed basis
signs. Write \[
 U_j=R_j+R_8\ (j\le6),\qquad U'_6=R_6+R_9,\qquad
 U_7=R_7+R_9=x_7-x_6,\qquad V_j=U_j+R_9\ (j\le6).
\] Every displayed \(U\) is a CG-positive ternary word of value 2; every
\(V\) is one of value 3, and \(V_6=x_6-x_5\) is always positive in
sorted order.

If a word \(p\) flips but \(p+k\) stays positive, a coefficient \(y\) on
\(p\) can be replaced using an available \(z\ge y\) from the positive
word \(k\): \[
 yp+zk=(z-y)(-p)+z(p+k).
\] This changes the lower bound by \(z-2y\). Several flipped basis rows
can share the same budget \(z_{\rm total}\) provided each allocation is
at least its own original weight. When \(U_j\) flips but \(V_j\) stays
positive, instead use \[
 y_jR_j+y_jR_8+zR_9=(z-y_j)(-U_j)+zV_j,\qquad z\ge y_j.
\] The gain is \(z-3y_j\), using only \(y_j\) of \(R_8\) and \(z\) of
\(R_9\). All unallocated positive terms are retained. Only these simple
substitutions and three integer identities below are required.

\subsection{Complete partition by flipped basis
rows}\label{complete-partition-by-flipped-basis-rows}

Any chamber of sign distance at most two is already covered. For three
flips, let \(m\) be the number of basis rows among them.

\begin{itemize}
\item
  \textbf{\(m=0\).} The original positive identity is unchanged.
\item
  \textbf{\(m=1\), basis \(R_j\), \(j\le5\).} If \(U_j\) remains
  positive, the single-flip proof works. If \(U_j\) flips but \(V_j\)
  stays positive, use the second substitution with \(z=57\), as in
  radius two. The only remaining flip set is \(\{R_j,U_j,V_j\}\).
  Rewrite \(y_jR_j+y_jR_8+y_jR_9=y_jV_j\). Now apply the first
  substitution with flipped word \(p=V_j\) and \(k=R_6\) for \(j=1,2\),
  \(k=R_7\) for \(j=3,4\), and \(k=R_8\) for \(j=5\), using that word's
  entire remaining budget. The new word \(V_j+k\) is ternary and
  CG-positive of value 4. None of the three flips is that word or the
  reserve \(k\). The resulting bounds for \(j=1,\ldots,5\) are
  \(172,172,178,170,169\). For \(j=6\), if either \(U_6\) or \(U'_6\)
  stays positive, use that companion and its reserve. Otherwise both are
  flipped, and the second substitution with \(U_6,V_6\) gives 173;
  \(V_6\) cannot flip in a sorted chamber. For \(j=7\), use the original
  \(U_7\) certificate regardless of additional flips.
\item
  \textbf{\(m=2\).} Start with the two-basis radius-two proof: rows
  \(R_1,\ldots,R_5\) use the \(R_8\) reserve, \(R_6,R_7\) use \(R_9\) if
  only one of them is flipped, and if both are flipped \(R_6\) uses
  \(R_8\) while \(R_7\) uses \(R_9\). The extra flip can affect this
  proof \textbf{only} if it is a designated companion; \(U_7\) cannot
  flip. Each such relevant case is covered as follows.

  \begin{itemize}
  \tightlist
  \item
    If a lower index \(j\le5\) has its \(U_j\) flipped and the other
    basis index \(i\) is at most 6, use the second substitution for
    \(j\) with \(z=57\); use the remaining \(48-y_j\) of \(R_8\) in the
    first substitution for \(i\) (so if \(i=6\), its companion is
    \(U_6\)). Its needed companion is different from \(U_j\), since
    adding the same \(R_8\) to two distinct \(R\)'s is injective.
  \item
    If \(j\le5\) has \(U_j\) flipped and the other basis is 7, use
    \(z=32\) for \(j\), leaving 25 of \(R_9\) for the first substitution
    for \(R_7\). This gives \(161+7-3y_j\ge162\) for \(j=1,2,3\). For
    \(j=4,5\) use the identities in the next section.
  \item
    If \(R_6\) and a lower basis row flip, and the extra change is
    \(U'_6\), use \(R_8\) for both substitutions (their required
    allocations total at most \(15+8<48\)); \(U_6\) remains positive. If
    \(R_6,R_7,U_6\) flip, use the last integer identity below.
  \end{itemize}

  This is the whole list of harmful companion changes: each of the 21
  pairs has two designated companions except that \(U_7\) is
  sorted-positive for the six pairs containing \(R_7\). Hence there are
  \(42-6=36\) exact obligations in this subfamily; an unrelated extra
  flip leaves the already verified radius-two dual untouched.
\item
  \textbf{\(m=3\).} No companion flips. Use the same reserve rule: lower
  indices \(1\) through \(5\) share \(R_8\); \(R_6\) or \(R_7\) alone
  among the two uses \(R_9\); when both occur, \(R_6\) shares \(R_8\)
  with the one lower row and \(R_7\) uses \(R_9\). Three lower rows have
  weight sum at most 14, so their common reserve gives gain at least
  \(48-2\cdot14\). If both higher rows occur, the shared lower reserve
  has weights summing to at most \(15+8=23\), and contributes gain at
  least 2, plus 7 from \(R_7\). The remaining mixed case uses disjoint
  reserves and has even larger gain. The 35 three-basis sets are all
  verified, with minimum resulting proved lower bound 170.
\end{itemize}

\subsection{Three explicit exceptional integer
duals}\label{three-explicit-exceptional-integer-duals}

For the flip sets \(\{R_j,R_7,U_j\}\) with \(j=4,5,6\), set
\(\epsilon=0\) if \(j=4\) and 1 otherwise, and let
\(Q=R_j+R_7+R_9+\epsilon R_8\). The table gives the new signed words,
all coefficients on the rows involved, and the overall lower bound. Keep
the original terms \((R_i,y_i)\) for \(i\notin\{j,7,8,9\}\).

\begin{longtable}[]{@{}
  >{\raggedright\arraybackslash}p{(\columnwidth - 12\tabcolsep) * \real{0.1154}}
  >{\raggedright\arraybackslash}p{(\columnwidth - 12\tabcolsep) * \real{0.1154}}
  >{\raggedleft\arraybackslash}p{(\columnwidth - 12\tabcolsep) * \real{0.1538}}
  >{\raggedleft\arraybackslash}p{(\columnwidth - 12\tabcolsep) * \real{0.1538}}
  >{\raggedleft\arraybackslash}p{(\columnwidth - 12\tabcolsep) * \real{0.1538}}
  >{\raggedleft\arraybackslash}p{(\columnwidth - 12\tabcolsep) * \real{0.1538}}
  >{\raggedleft\arraybackslash}p{(\columnwidth - 12\tabcolsep) * \real{0.1538}}@{}}
\toprule\noalign{}
\begin{minipage}[b]{\linewidth}\raggedright
\(j\)
\end{minipage} & \begin{minipage}[b]{\linewidth}\raggedright
CG-positive word \(Q\)
\end{minipage} & \begin{minipage}[b]{\linewidth}\raggedleft
on \(Q\)
\end{minipage} & \begin{minipage}[b]{\linewidth}\raggedleft
on \(-U_j\)
\end{minipage} & \begin{minipage}[b]{\linewidth}\raggedleft
on \(-R_7\)
\end{minipage} & \begin{minipage}[b]{\linewidth}\raggedleft
on \(R_8\)
\end{minipage} & \begin{minipage}[b]{\linewidth}\raggedleft
cap lower
\end{minipage} \\
\midrule\noalign{}
\endhead
\bottomrule\noalign{}
\endlastfoot
4 & ..-++.+-- & 57 & 53 & 32 & 101 & 270 \\
5 & \ldots-+.+.- & 57 & 49 & 32 & 40 & 201 \\
6 & \ldots.-.++- & 57 & 42 & 32 & 33 & 180 \\
\end{longtable}

The identities hold because the four replacements are \[
57Q+(57-y_j)(-U_j)+32(-R_7)
 +(105-y_j-57\epsilon)R_8
 = y_jR_j+25R_7+48R_8+57R_9.
\] Every \(Q\) is ternary and has positive CG value \(3+\epsilon\), so
it is not among the flips. All multipliers are nonnegative, completing
the proof.

\subsection{Exact replay}\label{exact-replay}

\begin{verbatim}
sha256sum -c SHA256SUMS
python3 check_radius3.py
\end{verbatim}

The utility checks all actual used words have ternary coefficients,
checks their CG orientation by integer dot products, verifies the model
sign for each term, and checks exact objective identity and nonnegative
multipliers in each reduced case, including the three exceptional
certificates. It checks six one-basis chain cases, 36 two-basis
companion cases and 35 three-basis cases. The written case partition
accounts for arbitrary other flips without an exhaustive search through
all sign patterns. The radius-two archive can separately be replayed as
documented there. No numerical optimization or Lean checking is claimed.

\section{Provenance and scope}\label{provenance-and-scope}

Correspondence:
\href{mailto:aafimalek2032@gmail.com}{\nolinkurl{aafimalek2032@gmail.com}}.
The proof, code, and writing in this research note were developed with
AI assistance.

See also the
\href{https://github.com/Aafimalek/novel-discoveries/tree/cfe880a69f3cf904c1c49dc531ecb18b28770964/nine-separated-subset-sums/verification/conway-guy-chambers-radius-two}{earlier
radius-two proof} and the unrelated
\href{https://github.com/Aafimalek/novel-discoveries/tree/76beffd491584923083c1c3d8f2abe637dde12a3/nine-separated-subset-sums/verification/lean-real-149}{Lean-checked
global lower bound of 149}. The present theorem concerns only the
bounded neighborhood of the given witness. It was checked using exact
integer arithmetic, not Lean.

\end{document}
